% Generated draft: requires mathematical and source review.
\documentclass[11pt]{article}
\usepackage[margin=25mm]{geometry}
\usepackage{fontspec}
\setmainfont{Latin Modern Roman}
\usepackage{amsmath,amssymb,mathtools,unicode-math}
\setmathfont{Latin Modern Math}
\usepackage{enumitem,needspace,xurl,hyperref}
\setlength{\parindent}{0pt}
\setlength{\parskip}{5pt}
\setlength{\emergencystretch}{3em}
\newcommand{\N}{\mathbb{N}}
\newcommand{\Z}{\mathbb{Z}}
\newcommand{\Q}{\mathbb{Q}}
\newcommand{\R}{\mathbb{R}}
\newcommand{\C}{\mathbb{C}}
\newcommand{\F}{\mathbb{F}}
\newcommand{\A}{\mathbb{A}}
\newcommand{\PP}{\mathbb{P}}
\newcommand{\E}{\mathbb{E}}
\newcommand{\eps}{\varepsilon}
\newcommand{\dd}{\,\mathrm{d}}
\begin{document}
\section*{UnsolvedMath ID: 6700037}
\textbf{Source problem code:} catalogue entry
\subsection{Definitions and mathematical statement}
Let $(X,g)$ be a Riemannian manifold with parallel Ricci tensor, $\nabla\operatorname{Ric}_g=0$, and positive Ricci tensor. Under the strict interpretation, positivity means $\operatorname{Ric}_g(v,v)>0$ for every nonzero tangent vector $v$. A local flat factor means a local Riemannian product factor with Euclidean metric.

The source conjectures that all such manifolds are area extremal, and those with no local flat factor are area rigid. It highlights symmetric spaces, Einstein spaces, and compact Lie groups with bi-invariant metrics as intended examples, subject to the actual curvature hypotheses. An Einstein metric has $\operatorname{Ric}_g=\lambda g$; a locally symmetric metric has parallel full curvature tensor. Neither class automatically has positive Ricci curvature, so the words ``all symmetric and all Einstein spaces'' cannot by themselves remove the positivity requirement.

For context, a common same-manifold area comparison uses a second metric $g'$ with
\[
|v\wedge w|_{g'}\geq|v\wedge w|_g
\]
for every tangent two-plane. One version of area extremality asks that $\operatorname{Sc}_{g'}\geq\operatorname{Sc}_g$ under this comparison force equality of scalar curvatures; a rigidity version also forces the metrics to agree up to the specified isometry relation. The source does not supply its definitions of area extremality/rigidity, which can also involve maps and degree hypotheses, so this context is not substituted for a verified exact convention.

The no-flat-factor qualifier is retained even though strictly positive Ricci curvature already excludes such a factor. This suggests a possible nonnegative-versus-positive convention that the extract does not resolve. Compactness, completeness, and the comparison class are also unstated.
\subsection{Short English statement}
Prove the proposed area extremality and rigidity of manifolds with positive parallel Ricci tensor, with no-flat-factor rigidity. Specify the area-comparison definition and clarify the positivity and example-class qualifications.
\subsection{Sources}
\begin{itemize}
\item Misha Gromov, \emph{101 Questions, Problems and Conjectures around Scalar Curvature} (2017), printed tag [?40], occurrence 39, PDF page 48, \url{https://www.ihes.fr/~gromov/wp-content/uploads/2018/08/101-problemsOct1-2017.pdf}; accessed 2026-07-29, PDF-text-normalized extraction. The provenance describes the public author PDF as incomplete.
\item The source cites a local extremality lemma in reference [39], without supplying its bibliography entry, and discusses trivial tangent bundles of compact bi-invariant Lie groups as an obstacle to available Dirac-index proofs.
\item The supplied summary mentions Brendle--Marques--Neves, Gromov, Dirac-operator/minimal-hypersurface methods, and spherical/projective/hyperbolic examples without full theorem citations or hypotheses.
\item Local registry \texttt{data/unsolvedmath/problems.json}, permanent UnsolvedMath ID $6700037$, source code catalogue entry, \emph{Scalar Curvature Question [?40]}. Dataset attribution: Ulam AI and the UnsolvedMath Contributors.
\end{itemize}
\end{document}
